% Simulazione d'esame: prova del 01/02/1970, esercizio 2 -- pistone con legge parabolica e poi lineare.
% Stessi numeri e figure dello svolgimento di riferimento:
% svolgimenti/1970-02-01_E02_pistone_parabolico (calcoli.py e figure/).
% Compilare con: bash ../compila.sh 1970-02-01_E02_pistone_parabolico  (A4 + orologio)

% Preambolo comune delle simulazioni d'esame (vedi GUIDA.md).
% Lo stesso sorgente si compila in due formati:
%   - A4 (svolgimento.pdf), da leggere su carta, PC o telefono;
%   - orologio (orologio.pdf): colonna stretta, letta dall'app dell'orologio
%     (garmin-latex), che taglia le pagine in strisce larghe quanto lo schermo.
% Il formato orologio si sceglie definendo \formatoorologio prima di \input
% (lo fa compila.sh). In quel formato il margine destro è volutamente largo:
% ciò che sporge oltre la colonna (equazioni lunghe, tabelle, figure) viene
% riconosciuto dall'app e mostrato rimpicciolito, con lo zoom.

\documentclass[11pt]{article}

\usepackage[T1]{fontenc}
\usepackage[utf8]{inputenc}
\usepackage[italian]{babel}
\usepackage{amsmath,amssymb}
\usepackage{array}
\usepackage{siunitx}
\usepackage{booktabs}
\usepackage{graphicx}
\usepackage{float}
\usepackage{caption}
\usepackage{enumitem}
\usepackage{ragged2e}
\usepackage{titlesec}
\usepackage{hyperref}

\hypersetup{hidelinks}
\setlength{\parindent}{0pt}
\setlength{\parskip}{4pt}
\allowdisplaybreaks
\sisetup{output-decimal-marker={,}, group-separator={\,}, per-mode=symbol,
         exponent-product=\cdot, range-phrase={\text{--}}}

% Notazione del formulario (come template/svolgimento.tex).
\newcommand{\Cp}{C^{+}}
\newcommand{\Cm}{C^{-}}
\newcommand{\Mr}{M_{r}}
\newcommand{\Ms}{M_{s}}

% Separatore fra relazioni nella stessa formula: \begin{gather*} A \ee B \end{gather*}.
% In A4 le relazioni stanno sulla stessa riga, nel formato orologio una sotto l'altra.
\ifdefined\formatoorologio \newcommand{\ee}{,\\} \else \newcommand{\ee}{,\qquad} \fi

% Titolo dell'esercizio: \esercizio{Prova del GG/MM/AAAA -- Esercizio n}{argomento}.
% Va anche nei metadati del PDF: l'app dell'orologio lo usa come titolo nell'elenco.
\newcommand{\esercizio}[2]{%
  \title{#1\\[2pt]\large #2}\author{}\date{}%
  \hypersetup{pdftitle={#1: #2}}}

\ifdefined\formatoorologio
  % Colonna di 4,6 cm: a 2,4 px/pt è larga 262 px sullo schermo da 390.
  \usepackage[paperwidth=12cm,paperheight=150cm,left=0.2cm,top=0.2cm,
              textwidth=4.6cm,textheight=149cm,nohead,nofoot]{geometry}
  \usepackage[scaled=1.0]{helvet}
  \usepackage{sfmath}
  \renewcommand{\familydefault}{\sfdefault}
  \pagestyle{empty}
  \makeatletter\let\ps@plain\ps@empty\makeatother   % anche la pagina del titolo
  % colonna dichiarata per l'app dell'orologio (deve coincidere con left e textwidth)
  \hypersetup{pdfkeywords={colonna=0.2cm+4.6cm}}
  \RaggedRight
  \setlength{\parskip}{5pt}
  \hfuzz=20cm          % ciò che sporge va nel margine destro, senza avvisi
  \hbadness=10000
  \titleformat*{\section}{\normalsize\bfseries\RaggedRight\hyphenpenalty=10000}
  \titleformat*{\subsection}{\normalsize\bfseries\RaggedRight}
  \setlist{leftmargin=*,itemsep=2pt}
  \captionsetup{font=small,justification=raggedright,singlelinecheck=false}
  % figure: larghe il doppio della colonna, così sporgono e l'app le mostra con lo zoom
  \newcommand{\figura}[3][0.7]{%
    \begin{figure}[H]\includegraphics[width=9cm]{#2}\caption{#3}\end{figure}}
\else
  \usepackage[a4paper,margin=2.5cm]{geometry}
  \newcommand{\figura}[3][0.7]{%
    \begin{figure}[H]\centering\includegraphics[width=#1\textwidth]{#2}\caption{#3}\end{figure}}
\fi

\graphicspath{{../../svolgimenti/1970-02-01_E02_pistone_parabolico/figure/}}

\newcommand{\uma}{u-a}
\newcommand{\upa}{u+a}

\esercizio{Prova del 01/02/1970 -- Esercizio 2}{Pistone con legge parabolica e poi lineare}

\begin{document}
\maketitle

\emph{Traccia.} Un pistone all'estremo sinistro di un condotto molto lungo si muove verso destra con legge $x=t^2$
fino a $t=1$ e poi con legge $x=2t-1$. Il gas è in quiete con $p=1$ e $\rho=\gamma$. Determinare l'evoluzione del
campo e la sua dipendenza dal peso molecolare del gas ($\gamma$ e $R$); descrivere il campo nei piani $x$--$t$ e
$p$--$u$. (La traccia coincide con B69-E02 del 25/01/2023; dello stesso tipo B49-E01 e B20-E01.)

\section*{Che cosa dicono i dati}
Dalla legge del pistone ricavo
\begin{gather*}
x_p=t^2 \ee u_p=2t \ee \ddot x_p=2\quad(0\le t\le1),\\
x_p=2t-1 \ee u_p=2 \ee \ddot x_p=0\quad(t>1).
\end{gather*}
Posizione e velocità sono continue in $t=1$, salta solo l'accelerazione; anche in $t=0$ la velocità è continua
($u_p(0)=0$) e l'accelerazione passa da 0 a 2.

Il dato $\rho=\gamma$ con $p=1$ fissa $a_0=\sqrt{\gamma p_0/\rho_0}=1$ \textbf{per qualunque gas}: le velocità sono
riferite ad $a_0$ e il pistone arriva a $u_p=2a_0$. La temperatura iniziale invece dipende dal gas,
$T_0=p_0/(\rho_0R)=1/(\gamma R)$. La traccia non fissa il gas: lavoro in forma letterale in $\gamma$, do i numeri per
l'aria ($\gamma=1{,}4$) e alla fine confronto elio ($\gamma=5/3$) e SF$_6$ ($\gamma\approx1{,}1$).

\section*{Modello}
Uso Eulero 1D instazionario, gas perfetto non viscoso e adiabatico: il condotto è a sezione costante e molto lungo,
quindi onde piane e nessuna riflessione dall'estremo destro. Finché non si forma l'urto il campo è un'onda semplice
omoentropica, perché lo stato iniziale è uniforme e il pistone, partendo da fermo, genera onde di una sola famiglia.
Dopo la formazione il flusso non è più omoentropico: per gli stati dopo che l'urto ha inglobato tutta la
compressione uso un pseudo-problema di Riemann fra gli stati esterni, che dà gli stati medi ma non la struttura
interna delle zone d'interazione; il regime asintotico lo ricavo dal pistone a velocità costante. Gli istanti
d'incontro dopo la formazione li stimo con caratteristiche rettilinee anche nelle zone d'interazione.

\section*{Che cosa mi aspetto}
Il pistone è a sinistra del gas e avanza: genera una \textbf{compressione $\upa$}, che spinge il gas verso destra,
nel verso in cui si propaga. Poiché $u_p(0)=0$ non c'è singolarità in $(0,0)$: la testa è la caratteristica $x=a_0t$,
su cui è discontinuo solo il gradiente, e la compressione è \textbf{non centrata}. Ogni caratteristica è una retta che
porta lo stato del pistone all'istante di emissione; quelle emesse più tardi hanno $u$ e $a$ maggiori, sono più
veloci e raggiungono le precedenti: la compressione si irripidisce e degenera in un urto per \textbf{coalescenza}.
L'urto nasce di intensità nulla e si rinforza assorbendo le caratteristiche che lo raggiungono da sinistra.

Dopo $t=1$ il pistone ha velocità costante: l'ultima caratteristica di compressione (la coda) parte da $(1,1)$ e fra
essa e il pistone c'è un plateau uniforme, lo stato 1. Il gas è stato compresso isentropicamente, mentre a parità di
velocità del gas un urto sostiene una pressione minore (la Hugoniot sta sotto l'isentropica). Quando l'urto avrà
inglobato tutta la compressione, mi aspetto quindi, come al centro di una compressione centrata, un \textbf{urto
$\upa$, un contatto e un'espansione $\uma$} (RCS). L'espansione torna sul pistone, si riflette con la stessa natura,
insegue l'urto e lo indebolisce; a tempi lunghi resta il campo del pistone impulsivo a velocità 2, con la memoria del
transitorio, uno strato entropico.

Gli stati che userò sono: 0 gas in quiete; 1 plateau sul pistone ($u=2$); 2 e 3 dietro l'espansione $\uma$ e dietro
l'urto dopo il pseudo-Riemann 1|0; 4 sul pistone dopo la riflessione; 5 (gas mai urtato sul pistone) e 6 (dietro
l'urto) nel regime asintotico. Nel piano $x$--$t$ avrò una zona a 0 famiglie (0, poi i plateau 1 e 4), una zona a una
famiglia (l'onda semplice, con le $\upa$ semirette, delimitata dalle semirette testa e coda finché l'urto non le
raggiunge) e, dopo l'urto, zone d'interazione a due famiglie vicino al pistone e dove l'espansione riflessa attraversa
il contatto e raggiunge l'urto.

\section*{L'onda semplice di compressione}
In tutta l'onda semplice $\upa$ l'invariante dell'altra famiglia è quello del gas in quiete, $a-\delta u=a_0=1$:
questa è la scorciatoia che mi dà ogni caratteristica in forma chiusa. Quella emessa all'istante $t_p\le1$ porta
\begin{gather*}
u=2t_p \ee a=1+2\delta t_p \ee \lambda=u+a=1+(\gamma+1)t_p,
\end{gather*}
quindi è la retta
\[
x=t_p^2+\left[1+(\gamma+1)t_p\right](t-t_p),
\]
con $p/p_0=(a/a_0)^{2\gamma/(\gamma-1)}$ e $\rho/\rho_0=(a/a_0)^{2/(\gamma-1)}$. La testa ($t_p=0$) è $x=t$; la coda
($t_p=1$) parte da $(1,1)$ con pendenza $2+\gamma$ (aria: $3{,}4$).

\section*{Dove si forma l'urto}
Due caratteristiche vicine si incontrano dove $\partial x/\partial t_p=0$ a $t$ fissato:
\[
2t_p+(\gamma+1)(t-t_p)-\left[1+(\gamma+1)t_p\right]=0
\]
\[
\Rightarrow\ t_{env}(t_p)=\frac{1+2\gamma t_p}{\gamma+1}.
\]
È l'inviluppo, la frontiera della soluzione a più valori. Cresce con $t_p$ (controllo: il primo punto di coalescenza
è sulla testa, $t_p=0$), quindi
\[
t_{BU}=x_{BU}=\frac{1}{\gamma+1}=\frac{a_0}{(1+\delta)\,\ddot x_p(0)}
\]
\[
\text{aria: } t_{BU}=x_{BU}=0{,}417.
\]
Controllo subito che l'urto nasca durante la fase parabolica e davanti al pistone: $t_{BU}=0{,}417<1$ e
$x_{BU}=0{,}417>x_p(t_{BU})=0{,}174$. Il pistone raggiungerebbe la testa $x=t$ solo in $t=1$, ma lì c'è già l'urto,
più veloce della testa.

\figura[0.78]{xt_formazione.pdf}{Piano $x$--$t$ fino a $t=1{,}8$ (aria): caratteristiche $\upa$ (tratto fisico
continuo, prolungamento a più valori punteggiato fino all'inviluppo), urto stimato da BU ad A, contatto ed
espansione $\uma$ del pseudo-Riemann; ombreggiato il gas urtato durante la crescita dell'urto.\label{fig:xtf}}

Nella fig.~\ref{fig:xtf} le caratteristiche si infittiscono vicino alla testa, dove l'inviluppo (tratteggio grigio)
ha il primo punto in BU. Senza l'urto la soluzione sarebbe a più valori fra testa e inviluppo: l'urto taglia questa
regione e assorbe una dopo l'altra le caratteristiche, fino alla coda, che raggiunge in A. La particella che l'urto
investe in BU (punteggiata) separa il gas mai urtato, una striscia sottile vicino al pistone, dal gas urtato da un
urto via via più forte.

\section*{Il plateau 1}
Dalla coda ($t_p=1$, $u_1=u_p=2$):
\begin{gather*}
a_1=a_0+\delta u_p=1+(\gamma-1)=\gamma \ee p_1=\gamma^{\frac{2\gamma}{\gamma-1}},\\
\frac{\rho_1}{\rho_0}=\gamma^{\frac{2}{\gamma-1}} \ee \frac{T_1}{T_0}=\gamma^2 \ee M_1=\frac{2}{\gamma}.
\end{gather*}
Aria: $a_1=1{,}4$, $p_1=10{,}5$, $\rho_1/\rho_0=5{,}38$, $T_1/T_0=1{,}96$, $M_1=1{,}43$. Il plateau è supersonico:
le $\uma$ al suo interno hanno pendenza $u_1-a_1=2-\gamma=0{,}6>0$, vanno verso destra nel laboratorio ma più
lentamente del pistone, che quindi le raggiunge.

Prima di proseguire controllo se potrei trattare l'urto come un'onda isentropica: con $M_s\approx2{,}8$, ben oltre
$\approx1{,}2$, no. Lo stato 1 sovrastima del 20\,\% la pressione che un urto sostiene alla stessa velocità, quindi
serve Rankine--Hugoniot.

\section*{L'urto cresce fino ad A}
Fra BU e A l'urto assorbe le caratteristiche a partire dalla testa. Ne stimo la traiettoria per punti: quando
raggiunge la caratteristica emessa in $t_p$, gli assegno la velocità dell'urto del problema di Riemann fra lo stato
portato da quella caratteristica e il gas in quiete, trascurando le onde che l'urto stesso riflette. La velocità
cresce da $w=a_0=1$ in BU al valore del pseudo-Riemann qui sotto. Controllo che lungo tutta la traiettoria la velocità
stia fra quella delle caratteristiche davanti ($1$) e dietro ($1+(\gamma+1)t_p$) e che l'urto resti prima
dell'inviluppo: è così. L'urto raggiunge la coda in
\begin{gather*}
\text{A}:\quad t_A\approx1{,}34 \ee x_A\approx2{,}16\quad(\text{aria}),
\end{gather*}
con il pistone in $x_p\approx1{,}68$: lì ha inglobato tutta la compressione.

\section*{In A: pseudo-Riemann fra 1 e 0}
Metto a contatto lo stato 1 (sinistra) e lo stato 0 (destra). Le curve nel piano $p$--$u$ sono
\[
u_L(p)=u_1-(p-p_1)\sqrt{\frac{A_1}{p+B_1}}\quad(p>p_1),
\]
\[
u_L(p)=u_1+\frac{a_1}{\delta}\left[1-\left(\frac{p}{p_1}\right)^{\frac{\gamma-1}{2\gamma}}\right]\quad(p\le p_1),
\]
\[
u_R(p)=u_0+(p-p_0)\sqrt{\frac{A_0}{p+B_0}}\quad(p>p_0),
\]
con $A_i=2/[(\gamma+1)\rho_i]$ e $B_i=(\gamma-1)p_i/(\gamma+1)$.

\textbf{Primo tentativo isentropico.}
\[
z=\frac{a_0}{a_1}\left(\frac{p_1}{p_0}\right)^{\frac{\gamma-1}{2\gamma}}=1,
\]
\[
u^{(0)}=\frac{z(a_1+\delta u_1)-(a_0-\delta u_0)}{\delta(1+z)}=2.
\]
Il tentativo isentropico mi restituisce lo stato 1, perché 1 e 0 stanno sulla stessa isentropica $\upa$: è la
seconda scorciatoia, il residuo misura direttamente la differenza fra isentropica e Hugoniot. Con $u^{(0)}=2$ l'onda
sinistra è nulla, $p_2^{(0)}=p_1=10{,}5$, mentre l'urto destro dà $p_3^{(0)}=8{,}73$. A parità di $u$ l'urto sostiene
meno pressione, quindi $p^*$ deve scendere sotto $p_1$ e $u^*$ salire sopra 2: a sinistra c'è un'espansione.
Aggiorno $u$ finché $|p_2-p_3|$ è trascurabile:
\begin{gather*}
u^*=u_2=u_3=2{,}11 \ee p^*=p_2=p_3=9{,}48\quad(\text{RCS}).
\end{gather*}
Controllo la natura delle onde: $u^*>u_1=2$, espansione a sinistra; $u^*>u_0=0$, urto a destra. Verificata.

Ricostruisco gli stati, con l'isentropica a sinistra e Rankine--Hugoniot a destra:
\begin{gather*}
a_2=a_1-\delta(u^*-u_1)=1{,}38 \ee \frac{\rho_2}{\rho_0}=4{,}98 \ee \frac{T_2}{T_0}=1{,}90,
\end{gather*}
\[
\frac{\rho_3}{\rho_0}=\frac{(\gamma+1)p^*+(\gamma-1)p_0}{(\gamma-1)p^*+(\gamma+1)p_0}=3{,}74,
\]
\begin{gather*}
a_3=1{,}59 \ee \frac{T_3}{T_0}=2{,}53.
\end{gather*}
Velocità dell'urto:
\[
M_{rel,0}=\frac{\gamma+1}{4}\frac{u^*}{a_0}+\sqrt{1+\left(\frac{\gamma+1}{4}\frac{u^*}{a_0}\right)^2},
\]
\[
w=u_0+a_0M_{rel,0}\ \Rightarrow\ w=M_s=2{,}87.
\]
Controllo l'urto: $M_s=2{,}87>1$ e $s_3/s_0=1{,}50>1$ ($\Delta s/R=1{,}01$). Verificato. L'espansione $\uma$ ha testa
$u_1-a_1=0{,}60$ e coda $u^*-a_2=0{,}73$; il contatto va a $u^*=2{,}11$.

\section*{L'espansione torna sul pistone: stato 4}
Con caratteristiche rettilinee la testa dell'espansione raggiunge il pistone in P ($t\approx1{,}69$, $x\approx2{,}37$),
la coda in $t\approx1{,}72$. Sul pistone c'è una sola caratteristica entrante (la $\upa$), quindi impongo una sola
condizione, $u=u_p=2$; lungo la $\uma$ arriva l'invariante $a-\delta u$ dello stato 2:
\begin{gather*}
a_4-\delta u_p=a_2-\delta u^* \ \Rightarrow\ a_4=a_2-\delta(u^*-u_p)=1{,}36,
\end{gather*}
\[
p_4=p_2\left(\frac{a_4}{a_2}\right)^{\frac{2\gamma}{\gamma-1}}=8{,}50.
\]
Controllo: $a_4<a_2$, la riflessa è un'espansione $\upa$, come deve essere sulla parete mobile (stessa natura); e
$a>0$, $p>0$ ovunque. La sua testa, con pendenza $u^*+a_2=3{,}48$, attraversa il contatto in C ($t\approx2{,}06$,
$x\approx3{,}67$) e, con pendenza $u^*+a_3=3{,}70$, raggiunge l'urto in S ($t\approx2{,}73$, $x\approx6{,}1$): da lì
l'urto rallenta.

Al contatto l'espansione passa dal gas 2 ($\rho_2a_2=9{,}62$) al gas 3 ($\rho_3a_3=8{,}33$), verso impedenza minore:
oltre all'espansione trasmessa nasce una debole compressione riflessa, che torna sul pistone e riporta la pressione
da $p_4=8{,}50$ verso il valore asintotico.

\figura[0.8]{xt_campo.pdf}{Piano $x$--$t$ dopo A (aria), con caratteristiche rettilinee: espansione $\uma$ da A al
pistone (P), testa dell'espansione riflessa $\upa$ attraverso il contatto (C) fino all'urto (S); dopo S l'urto
rallenta verso $w_\infty$ (raccordo qualitativo).\label{fig:xtc}}

Nella fig.~\ref{fig:xtc} pistone (2), contatto (2,11) e urto (2,87) hanno pendenze vicine: i gas 1 e 2 occupano zone
sottili fra pistone e contatto, e l'espansione $\uma$, pur andando verso destra, è raggiunta dal pistone in P. Dopo S
la traiettoria dell'urto si incurva verso la pendenza asintotica 2,76, mentre fra pistone e urto si stabiliscono gli
stati 5 e 6.

\section*{A tempi lunghi: stati 5 e 6}
Quando tutte le onde del transitorio hanno raggiunto l'urto, il gas dietro l'urto ha la velocità del pistone,
$u_6=u_p=2$: è la terza scorciatoia, l'urto è quello del pistone impulsivo. Da massa e quantità di moto nel
riferimento dell'urto
\begin{gather*}
\rho_0w=\rho_6(w-u_p) \ee p_0+\rho_0w^2=p_6+\rho_6(w-u_p)^2,
\end{gather*}
\[
\Rightarrow\ p_6=p_0+\rho_0\,w\,u_p,
\]
e con l'energia ottengo
\[
w^2-(1+\delta)u_pw-a_0^2=0,
\]
\[
w_\infty=\frac{(1+\delta)u_p+\sqrt{(1+\delta)^2u_p^2+4a_0^2}}{2}=2{,}76,
\]
\begin{gather*}
p_6=p_5=8{,}73 \ee \frac{\rho_6}{\rho_0}=\frac{w_\infty}{w_\infty-u_p}=3{,}62 \ee \frac{T_6}{T_0}=2{,}41.
\end{gather*}
Il gas a contatto con il pistone (5) non è mai stato urtato: ha l'entropia iniziale, $u_5=2$ e $p_5=p_6$:
\[
a_5=a_0\left(\frac{p_5}{p_0}\right)^{\frac{\gamma-1}{2\gamma}}=1{,}36,
\]
\begin{gather*}
\frac{\rho_5}{\rho_0}=4{,}70 \ee \frac{T_5}{T_0}=1{,}86.
\end{gather*}
Controllo: $p_5=p_6$, $u_5=u_6=u_p$, $p_6-p_0=\rho_0w_\infty u_p$; urto con $M_s=2{,}76>1$ e $s_6/s_0=1{,}44>1$
($\Delta s/R=0{,}91$, meno dell'urto in A). Verificato. Fra 5 e 6 c'è lo strato entropico: $p$ e $u$ uniformi,
densità e temperatura variabili.

\section*{Il piano $p$--$u$}
\figura[1.0]{pu.pdf}{Piano $p$--$u$ (aria), scala lineare e logaritmica: isentropica $\upa$ e Hugoniot $\upa$ da 0,
espansione $\uma$ da 1, espansione $\upa$ da 2 (riflessione sul pistone), retta $u=u_p=2$.\label{fig:pu}}

\figura[0.62]{pu_zoom.pdf}{Piano $p$--$u$ vicino a $u=2$ (aria): residuo del primo tentativo isentropico e stati del
transitorio.\label{fig:puz}}

Nella fig.~\ref{fig:pu} la compressione porta da 0 a 1 lungo l'isentropica $\upa$; la Hugoniot da 0 ha la stessa
tangente in 0 e sta sotto l'isentropica per $u>0$: le due coincidono finché la compressione è debole e si separano
sopra $u\approx1$. In scala lineare leggo le distanze in pressione vicino a $u=2$; in scala logaritmica la curvatura
è attenuata e si vede che le curve partono tangenti da 0. Nello zoom (fig.~\ref{fig:puz}) il segmento verticale fra 1
e la Hugoniot in $u=2$ è il residuo del primo tentativo; 2=3 è l'intersezione dell'espansione $\uma$ da 1 con la
Hugoniot; 4 sta sulla retta $u=2$, sotto la Hugoniot; il regime asintotico 5=6 sta sulla Hugoniot in $u=2$. Gli stati
5 e 6 coincidono in $p$--$u$ ma hanno impedenze diverse ($\rho_5a_5=8{,}97$, $\rho_6a_6=7{,}88$): il contatto non si
vede in questo piano.

\section*{Dipendenza dal gas}
In forma letterale $t_{BU}=1/(\gamma+1)$, $a_1=\gamma$, $p_1=\gamma^{2\gamma/(\gamma-1)}$ e $w_\infty$ (con
$1+\delta=(\gamma+1)/2$) dipendono solo da $\gamma$. Per $\gamma\to1$: $t_{BU}\to1/2$, $p_1\to e^2\approx7{,}39$,
$w_\infty\to1+\sqrt2\approx2{,}41$. Ripetendo il calcolo:

\begin{table}[H]
\centering\small
\begin{tabular}{lccc}
\toprule
 & He & aria & SF$_6$ \\
\midrule
$\gamma$ & 5/3 & 1,4 & 1,1 \\
$t_{BU}$ & 0,375 & 0,417 & 0,476 \\
$p_1$ & 12,9 & 10,5 & 8,14 \\
$M_1$ & 1,20 & 1,43 & 1,82 \\
$p^*$ & 11,8 & 9,48 & 7,14 \\
$u^*$ & 2,09 & 2,11 & 2,13 \\
$w$ (A) & 3,10 & 2,87 & 2,62 \\
$w_\infty$ & 3,00 & 2,76 & 2,50 \\
$p_6$ & 11,0 & 8,73 & 6,50 \\
$\rho_6/\rho_0$ & 3,00 & 3,62 & 5,00 \\
$T_6/T_0$ & 3,67 & 2,41 & 1,30 \\
$T_5/T_0$ & 2,61 & 1,86 & 1,19 \\
\bottomrule
\end{tabular}
\caption{Dipendenza dal gas ($p^*$, $u^*$, $w$ in A stime del modello; valori asintotici esatti).}
\end{table}

\figura[1.0]{gas.pdf}{Dipendenza da $\gamma$: istante di formazione, velocità dell'urto (pseudo-Riemann e
asintotica), pressioni dei plateau; le rette punteggiate indicano SF$_6$, aria ed elio.\label{fig:gas}}

Al crescere di $\gamma$ il blow-up è più precoce, l'urto più veloce e le pressioni più alte
(fig.~\ref{fig:gas}); la differenza fra la velocità del pseudo-Riemann e quella asintotica resta di circa
$0{,}1$--$0{,}12\,a_0$ per tutti i gas.

$R$ non compare: con i dati assegnati le equazioni di Eulero contengono solo $\gamma$, e $R$ entra solo nella
temperatura, $T=p/(\rho R)$. I rapporti $T/T_0$ sono indipendenti da $R$; cambia la temperatura assoluta,
$T_0=1/(\gamma R)$: a pari $p_0$ e $\rho_0$ un gas pesante (piccolo $R$) è più caldo.

\section*{Risposta}
Il pistone genera una compressione $\upa$ non centrata che degenera in un urto sulla testa in
$t_{BU}=x_{BU}=1/(\gamma+1)$ (aria $0{,}417$). Dietro la compressione c'è il plateau isentropico 1 ($u=2$,
$a=\gamma$, $p=\gamma^{2\gamma/(\gamma-1)}$, aria $10{,}5$). L'urto si rinforza fino ad A ($t\approx1{,}34$), dove
il pseudo-Riemann 1|0 dà RCS con $w=2{,}87$, $p^*=9{,}48$, $u^*=2{,}11$; l'espansione $\uma$ si riflette sul pistone
($p_4=8{,}50$), indebolisce l'urto e il campo tende al pistone impulsivo: $w_\infty=2{,}76$, $p=8{,}73$, con uno
strato entropico fra il gas sul pistone (5) e quello urtato (6). Gli andamenti con $\gamma$ sono nella tabella.

\section*{Discussione}
\textbf{Peso molecolare.} Un gas leggero monoatomico ha $\gamma$ alto, uno pesante poliatomico $\gamma$ vicino a 1. A
pari $a_0$, con $\gamma$ alto la velocità caratteristica $u+a=a_0+(1+\delta)u$ cresce di più con $u$: la non linearità
è più forte, l'urto si forma prima, è più veloce e la pressione è più alta (elio: $w_\infty=3$, $p_6=11$). Con
$\gamma\to1$ il gas è più comprimibile: urto più lento e pressione più bassa, ma densità dietro l'urto maggiore
(SF$_6$: $\rho_6/\rho_0=5$ contro 3 per l'elio) e riscaldamento minore ($T_6/T_0=1{,}30$ contro $3{,}67$). $R$ cambia
solo la scala delle temperature. Se si confrontassero gas reali alla stessa temperatura, il dato in unità di $a_0$
nasconderebbe il peso molecolare: $a_0=\sqrt{\gamma RT_0}$ a 293\,K è circa 1000\,m/s per l'elio e 135\,m/s per
SF$_6$, quindi la stessa legge del pistone in m/s sarebbe una compressione debole nell'elio e fortissima in SF$_6$.

\textbf{Isentropica contro urto.} Il pistone ``comprime meglio'' di un urto: lo stato 1 ha pressione maggiore di
quella che un urto sostiene alla stessa velocità del gas. Per questo, quando l'urto ha inglobato la compressione,
nasce l'espansione $\uma$ e l'urto è per un tratto più forte di quello asintotico ($2{,}87$ contro $2{,}76$);
l'espansione riflessa dal pistone lo riporta al valore del pistone impulsivo. È il meccanismo del problema di
Friedrichs, con la differenza che qui il pistone non si ferma.

\textbf{Strato entropico.} L'intensità dell'urto cresce da zero in BU fino al massimo vicino ad A e poi cala verso il
valore asintotico. Le particelle portano l'entropia ricevuta: fra il gas mai urtato vicino al pistone ($s_0$) e il gas
6 l'entropia prima cresce fino a $\approx s_3$ e poi scende a $s_6$. Lo strato è continuo, non c'è un vero contatto:
quello del pseudo-Riemann lo rappresenta in modo concentrato. Nel modello il gas fra pistone e contatto (1, 2, 4) ha
l'entropia iniziale; in realtà la maggior parte è stata urtata durante la crescita dell'urto, e il gas a $s_0$ è solo
una striscia sottile vicino al pistone. Gli stati 2, 3, 4 sono quindi indicativi, mentre il regime asintotico (urto,
stato 6, pressione) è esatto, perché dipende solo dalla velocità del pistone; anche gli istanti P, C, S sono stime,
perché nelle zone d'interazione le caratteristiche sono curve.

\textbf{Pressione sul pistone.} Cresce isentropicamente, $p_p=(1+2\delta t)^{2\gamma/(\gamma-1)}$, fino a
$p_1=10{,}5$ in $t=1$; resta costante finché non arrivano le onde riflesse dall'urto (nel modello l'espansione da A,
in $t\approx1{,}7$), scende a $\approx8{,}5$ e risale lentamente verso $8{,}73$ per le deboli compressioni riflesse
dal contatto.

\textbf{Limiti del modello.} Il modello 1D non viscoso trascura strato limite e scambio termico con le pareti; con
$T_6/T_0\approx2{,}4$ per l'aria a temperatura ambiente si arriva a circa 700\,K, dove il gas perfetto
caloricamente è ancora ragionevole.

\end{document}
